\subsection{Averaging Across Independent Sequences and Semantic Preservation}
\label{sec:sequence_level_averaging}

A central goal of this chapter is to isolate shared, grammar-like structure in
gradient matrices while preserving the model’s ability to learn semantic features.
This requires careful attention to \emph{where} averaging is performed and
\emph{what statistical independence assumptions} are invoked.

\paragraph{Token- and sequence-level structure.}
We index tokens by a pair $(j,t)$, where $j$ denotes the sequence index and $t$
denotes the token position within that sequence.
In typical language-model training, tokens within the same sequence share entities,
topics, and compositional meaning, and are therefore \emph{semantically correlated}.
In contrast, different sequences in a batch are drawn independently and may be
treated as approximately \emph{semantically independent}.

We decompose per-token activations and backpropagated errors as
\[
x^{(j,t)} = x_{\mathrm{gram}} + x_{\mathrm{sem}}^{(j,t)},
\qquad
\delta^{(j,t)} = \delta_{\mathrm{gram}} + \delta_{\mathrm{sem}}^{(j,t)},
\]
where the grammar components are shared across sequences, while the semantic
components vary across sequences and positions.

\paragraph{Batch gradient and semantic cancellation.}
The gradient of a linear projection with respect to a batch of tokens is
\[
G
=
\sum_{j} \sum_{t \in \text{seq } j}
\delta^{(j,t)} \bigl(x^{(j,t)}\bigr)^\top
\;\in\;
\mathbb{R}^{d\times d}.
\]
In this expression, the summation over tokens constitutes an implicit averaging
over the batch.
When this summation is taken over tokens from many \emph{independent sequences},
semantic contributions fluctuate in sign and direction and cancel in expectation,
while grammar-related contributions reinforce coherently.
This yields the decomposition
\[
G = A + R,
\]
where $A$ captures shared, grammar-like structure and $R$ captures residual
semantic variability, as formalized in Assumption~1.

\paragraph{Why averaging must be across sequences.}
It is crucial that this averaging be performed across tokens drawn from
\emph{different sequences}.
If averaging were instead restricted to tokens within a single sequence, semantic
signals would remain strongly correlated and would not cancel.
Such within-sequence averaging would suppress meaningful semantic structure and
interfere with learning.
In contrast, averaging across sequences removes only \emph{independent} semantic
variation, while preserving within-sequence correlations that are essential for
learning semantic and compositional features.

\paragraph{Averaging as a probe, not an update.}
The averaging operations introduced in this chapter (e.g., $G\mathbf{1}$ or
$G\Omega$) are used solely as \emph{diagnostic probes} to estimate the dominant
shared subspace of the gradient matrix.
They do not replace the full gradient in the parameter update.
All semantic information remains present in the original gradient and continues
to drive learning through stochastic gradient descent.
Averaging is therefore used to \emph{identify and regulate} dominant shared
directions, not to eliminate semantic gradients.

\paragraph{Implications for low-rank approximation.}
Under the assumption of semantic independence across sequences, the dominant
grammar component $A$ induces dense, coherent marginals that survive simple
aggregation, while the residual $R$ cancels under averaging.
This justifies the use of simple column averaging (and its randomized extensions)
to approximate the spiky, low-dimensional structure of the agreement matrix
without explicit singular value decomposition.
At the same time, because averaging is restricted to independent sequences,
semantic features remain learnable through standard optimization dynamics.

\paragraph{Summary.}
In summary, averaging should be understood as operating across independent
semantic units (sequences), not across correlated tokens within a sequence.
This distinction is essential: it allows shared grammatical structure to be
isolated without suppressing the semantic variability that drives representation
learning.